\documentclass{beamer}
\usepackage[net,green]{nesped}
\usepackage[brazilian]{babel}

% --- teste de fontes: Fraunces (títulos) + IBM Plex Sans (corpo) + JetBrains Mono (código) ---
% Fraunces restrito: SOFT=0, WONK=0 (sem o eixo "quirky"), instâncias estáticas em fonts/.
\setsansfont{IBM Plex Sans}[
    Path=fonts/,
    UprightFont=IBMPlexSans-Regular.ttf,
    BoldFont=IBMPlexSans-Bold.ttf,
    ItalicFont=IBMPlexSans-Italic.ttf,
]
\newfontfamily\titlefont{Fraunces SemiBold}[
    Path=fonts/,
    UprightFont=Fraunces-SemiBold.ttf,
    BoldFont=Fraunces-Bold.ttf,
]
\newfontfamily\codefont{JetBrains Mono}[
    Path=fonts/,
    UprightFont=JetBrainsMono-Regular.ttf,
    BoldFont=JetBrainsMono-Bold.ttf,
]
\setbeamerfont{title}{family=\titlefont, series=\bfseries}
\setbeamerfont{frametitle}{family=\titlefont, series=\bfseries}
\setmonofont{JetBrains Mono}[Path=fonts/, UprightFont=JetBrainsMono-Regular.ttf, BoldFont=JetBrainsMono-Bold.ttf]

\title[Teste]{Teste de fontes: Fraunces + IBM Plex Sans + JetBrains Mono}
\subtitle{Título com acentuação: relação, condição, não-inferioridade}
\author{Vitor Hugo}
\institute[UFV]{Universidade Federal de Viçosa --- PPGCC}
\date{28 de agosto de 2026}

\begin{document}

\titleframe{
    \titlelogo{img/logo-nesped.png}
    \titlelogo{img/logo-ufv.png}
}

\begin{frame}{Corpo de texto e acentuação}
    \begin{itemize}
        \item A representação em nível de check-in é o fator dominante.
        \item Não-inferioridade, margem de dois pontos, condição pré-registrada.
        \item Números: 0.19, 63.18\%, $p=0.011$, Acc@10.
    \end{itemize}
\end{frame}

\begin{frame}{Tabela com números em JetBrains Mono}
    \begin{table}
        \centering
        {\codefont
        \begin{tabular}{lrrr}
            Estado & Dedicado & Conjunto & $\Delta$ \\\hline
            Alabama & 63.18 & 63.44 & +0.26 \\
            Florida & 37.36 & 37.55 & +0.19 \\
            Texas   & 65.94 & 67.15 & +1.21 \\
        \end{tabular}}
    \end{table}
\end{frame}

\begin{frame}[fragile]{Código em JetBrains Mono}
    \begin{lstlisting}[language=Python]
def geom_simple(cat_f1, reg_acc10):
    return (cat_f1 * reg_acc10) ** 0.5
    \end{lstlisting}
\end{frame}

\end{document}
